% !TeX encoding = UTF-8
% !TeX spellcheck = es_MX

%Clase de documento, configuración de idioma, fuentes
\documentclass[a1paper,portrait,fontscale=0.40]{baposter}
\usepackage[spanish,es-tabla]{babel}
\usepackage[utf8]{inputenc}
\usepackage{arev}
\usepackage[T1]{fontenc}

%Matemáticas, símbolos
\usepackage{amsmath}
\usepackage{amssymb}

%Tipografía
\usepackage[auto]{microtype}

%Inclusión de imágenes, tablas, código fuente
\usepackage{graphicx}
\usepackage{booktabs}
\usepackage{multirow}
\usepackage{array}
\usepackage{paralist}
\usepackage{pifont}

%Colores y TikZ
\usepackage{xcolor}
\usepackage{colortbl}
\usepackage{tikz}
\usetikzlibrary{arrows.meta,positioning,shapes.geometric,backgrounds,fit,calc}

%Búsqueda de gráficos (las figuras/logos se reutilizan desde docs/poster/)
\graphicspath{{../sem_tesis4/figures/}{../../reports/Tesis/logo/}{../../reports/Tesis/figures/}{../../reports/}{bilder/}}

%Paleta del proyecto (de sem_tesis4 / poster navy-cream, con acento lila)
\definecolor{cream}{HTML}{F2EEE3}
\definecolor{paper}{HTML}{FBFAF6}
\definecolor{navy}{HTML}{2F3E4E}
\definecolor{accentA}{HTML}{7E57C2}
\definecolor{accentB}{HTML}{EF6C00}
\definecolor{accentC}{HTML}{2E7D32}
\definecolor{accentD}{HTML}{C62828}
\definecolor{rule}{HTML}{B5AFA1}

%Definición de colores (degradado lila en encabezados: blanco -> lila, marco negro)
\definecolor{standardfontcolor}{HTML}{2F3E4E}
\definecolor{bordercol}{HTML}{000000}
\definecolor{headercol1}{HTML}{FFFFFF}
\definecolor{headercol2}{HTML}{7E57C2}
\definecolor{headerfontcol}{HTML}{000000}
\definecolor{boxcolor}{HTML}{FBFAF6}

%Macro auxiliar para resaltar
\newcommand{\hl}[1]{{\color{accentD}\textbf{#1}}}

%Colores para los perfiles de aspectos (azul = Pueblo A, rosa = Pueblo B)
\definecolor{azulperfil}{HTML}{3B3FA8}
\definecolor{rosaperfil}{HTML}{C42B86}

%Diagrama de radar/red para los perfiles de aspectos de los pueblos.
% [#1] Opacidad del relleno (opcional, default 0.85)
%  #2  Título,  #3 Color,
%  #4..#8 Valores (Gastronomía, Naturaleza, Cultura, Tranquilidad, Precio) en escala 0-5.
\newcommand{\perfilradar}[8][0.85]{%
	\begin{tikzpicture}[every node/.style={font=\footnotesize}]
		\def\maxR{2.6}
		% Pentágonos concéntricos (rejilla)
		\foreach \r in {1,...,5}{
				\draw[navy!90,line width=0.4pt]
				(90:{\r/5*\maxR}) -- (18:{\r/5*\maxR}) -- (-54:{\r/5*\maxR}) --
				(-126:{\r/5*\maxR}) -- (162:{\r/5*\maxR}) -- cycle;
			}
		% Ejes
		\foreach \a in {90,18,-54,-126,162}{
				\draw[navy!35,line width=0.4pt] (0,0) -- (\a:\maxR);
			}
		% Polígono de datos con degradado radial (claro al centro -> intenso en el borde)
		\begin{scope}
			\clip (90:{#4/5*\maxR}) -- (18:{#5/5*\maxR}) -- (-54:{#6/5*\maxR}) --
			(-126:{#7/5*\maxR}) -- (162:{#8/5*\maxR}) -- cycle;
			\shade[inner color=#3!12,outer color=#3!85,opacity=#1] (0,0) circle (\maxR);
		\end{scope}
		\draw[draw=#3,line width=1.3pt,opacity=#1]
		(90:{#4/5*\maxR}) -- (18:{#5/5*\maxR}) -- (-54:{#6/5*\maxR}) --
		(-126:{#7/5*\maxR}) -- (162:{#8/5*\maxR}) -- cycle;
		% Círculo central con el promedio
		\fill[paper] (0,0) circle (0.74cm);
		\draw[navy!35,line width=0.4pt] (0,0) circle (0.74cm);
		\node[align=center] at (0,0) {Promedio:\\4.5};
		% Etiquetas de los ejes (radio mayor para separar del vértice)
		\pgfmathsetmacro{\labelR}{\maxR+0.14}
		\node[above]       at (90:\labelR)   {Gastronom\'ia};
		\node[right]       at (18:\labelR)   {Ambiente};
		\node[below right] at (-54:\labelR)  {Servicio};
		\node[below left]  at (-126:\labelR) {Infrastructura};
		\node[left]        at (162:\labelR)  {Precio};
		% Título
		\node[font=\normalsize\bfseries,text=#3] at (90:{\maxR+1.05}) {#2};
	\end{tikzpicture}%
}
